\documentclass[12pt]{article}

% Layout.
\usepackage[top=1.2in, bottom=0.75in, left=1in, right=1in, headheight=1.0in, headsep=0pt]{geometry}

% Fonts.
\usepackage{mathptmx}
\usepackage[scaled=0.86]{helvet}
\renewcommand{\emph}[1]{\textsf{\textbf{#1}}}

% TiKZ.
\usepackage{tikz, pgfplots}
\usetikzlibrary{calc}
\pgfplotsset{my style/.append style={axis x line=middle, axis y line=middle, xlabel={$x$}, ylabel={$y$}}}
\pgfplotsset{compat=1.16}

% Misc packages.
\usepackage{amsmath,amssymb,latexsym,bm,array,multicol,enumitem}
\usepackage{graphicx}
\usepackage{xcolor}
\usepackage[most]{tcolorbox}
\usepgfplotslibrary{fillbetween}
\definecolor{solbg}{RGB}{225,240,252}
\definecolor{solframe}{RGB}{70,130,180}
\newtcolorbox{solution}[1][]{%
  breakable, enhanced,
  colback=solbg, colframe=solframe, boxrule=0.7pt, arc=3pt,
  left=6pt, right=6pt, top=5pt, bottom=5pt,
  before skip=6pt, after skip=6pt,
  fonttitle=\sffamily\bfseries\small, coltitle=white, title={Solution},
  #1}

% Commands to set various header/footer components.
\makeatletter
\def\doctitle#1{\gdef\@doctitle{#1}}
\doctitle{Use {\tt\textbackslash doctitle\{MY LABEL\}}.}
\def\docdate#1{\gdef\@docdate{#1}}
\docdate{Use {\tt\textbackslash docdate\{MY DATE\}}.}
\def\doccourse#1{\gdef\@doccourse{#1}}
\let\@doccourse\@empty
\def\docscoring#1{\gdef\@docscoring{#1}}
\let\@docscoring\@empty
\def\docversion#1{\gdef\@docversion{#1}}
\let\@docversion\@empty
\makeatother

% Headers and footers layout.
\makeatletter
\usepackage{fancyhdr}
\pagestyle{fancy}
\fancyhf{} % Clears all headers/footers.
\lhead{\emph{\@doctitle\hfill\@docdate} \medskip
\ifnum \value{page} > 1\relax\else\\
\emph{Name: \rule{3.5in}{1pt}\hfill \@docscoring}
\fi}

\rfoot{\emph{\@docversion}}
\lfoot{\emph{\@doccourse}}
\cfoot{\emph{\thepage}}
\renewcommand{\headrulewidth}{0pt}%
\makeatother

% Paragraph spacing
\parindent 0pt
\parskip 6pt plus 1pt

% A problem is a section-like command. Use \problem{5} for a problem worth 5 points.
\newcounter{probcount}
\newcounter{subprobcount}
\setcounter{probcount}{0}

\newcommand{\problem}[1]{%
\par
\addvspace{4pt}%
\setcounter{subprobcount}{0}%
\stepcounter{probcount}%
\makebox[0pt][r]{\emph{\arabic{probcount}.}\hskip1ex}\emph{[#1 points]}\hskip1ex}

\newcommand{\thesubproblem}{\emph{\alph{subprobcount}.}}

% Subproblems are an enumerate-like environment with a consistent
% numbering scheme.  Use \begin{subproblems}\item...\item...\end{subproblems}
\newenvironment{subproblems}{%
\begin{enumerate}[itemindent=8pt,leftmargin=0pt]%
\setcounter{enumi}{\value{subprobcount}}%
\renewcommand{\labelenumi}{\emph{\textsl{\alph{enumi})}} \,}%
}%
{%
\setcounter{subprobcount}{\value{enumi}}%
\end{enumerate}%
}

% Blanks for answers in normal and math mode.
\newcommand{\blank}[1]{\rule{#1}{0.75pt}}
\newcommand{\mblank}[1]{\underline{\hspace{#1}}}
\def\emptybox(#1,#2){\framebox{\parbox[c][#2]{#1}{\rule{0pt}{0pt}}}}

% Misc.
\renewcommand{\d}{\displaystyle}
\newcommand{\ds}{\displaystyle}
\newcommand{\threeopts}{{\small \hspace{-6mm} $\begin{matrix} \text{\textsc{converges}} \\ \text{\textsc{absolutely}} \end{matrix}$ \qquad\qquad $\begin{matrix} \text{\textsc{converges}} \\ \text{\textsc{conditionally}} \end{matrix}$ \qquad\qquad \textsc{diverges}} \bigskip}


\newcommand{\ba}{\mathbf{a}}
\newcommand{\bb}{\mathbf{b}}
\newcommand{\bc}{\mathbf{c}}
\newcommand{\bi}{\mathbf{i}}
\newcommand{\bj}{\mathbf{j}}
\newcommand{\bk}{\mathbf{k}}
\newcommand{\bn}{\mathbf{n}}
\newcommand{\bu}{\mathbf{u}}
\newcommand{\bv}{\mathbf{v}}
\newcommand{\bw}{\mathbf{w}}
\newcommand{\br}{\mathbf{r}}
\newcommand{\bT}{\mathbf{T}}
\newcommand{\bN}{\mathbf{N}}
\newcommand{\bB}{\mathbf{B}}

\newcommand{\ip}[2]{\mathrm{\left<#1,#2\right>}}
\newcommand{\vv}[2]{\mathrm{\left<#1,#2\right>}}
\newcommand{\vvv}[3]{\mathrm{\left<#1,#2,#3\right>}}


\doctitle{Math 253: Quiz 5 Solutions}
\docdate{15 October, 2026}
\doccourse{}
\docversion{Solutions}
\docscoring{\fbox{{\LARGE \strut}\hspace{0.8in} / 25}}

\begin{document}
30 minutes maximum.  No aids (book, calculator, etc.) are permitted.  Show all work and use proper notation for full credit.  Answers should be in reasonably-simplified form.  25 points possible.

\begin{enumerate}

\item Consider the function $f(x,y) = x^2 + xy + y^2 - 3x$.
  \begin{enumerate}
  \item (4 points) Find all critical points of $f$.
    \begin{solution}
      \[
        f_x = 2x + y - 3 = 0, \qquad f_y = x + 2y = 0.
      \]
      The second equation gives $x = -2y$; substituting into the first,
      $-4y + y - 3 = 0$, so $y = -1$ and $x = 2$.
      \[
        \boxed{(2,\,-1)}
      \]
    \end{solution}
    \vfill

  \item (3 points) Classify the critical points using $\ds D(x_0,y_0) =
      f_{xx}(x_0,y_0)f_{yy}(x_0,y_0) - \left( f_{xy}(x_0,y_0)
      \right)^2$.
    \begin{solution}
      $f_{xx} = 2$, $f_{yy} = 2$, $f_{xy} = 1$, so
      \[
        D(2,-1) = (2)(2) - 1^2 = 3 > 0 \quad\text{and}\quad f_{xx}(2,-1) = 2 > 0.
      \]
      \[
        \boxed{f(2,-1) \text{ is a local minimum}}
      \]
    \end{solution}
      \vfill
    \end{enumerate}

    \newpage
    
\item (6 points) Find an equation for the tangent plane to the surface $-4x^3 + 12y^2 + 4z = -4$ at
  the point $(2,1,4)$.
  \begin{solution}
    Solving for $z$: $4z = 4x^3 - 12y^2 - 4$, so $z = x^3 - 3y^2 - 1$.
    \[
      \frac{\partial z}{\partial x} = 3x^2, \quad \frac{\partial z}{\partial x}\Big|_{(2,1)} = 12;
      \qquad
      \frac{\partial z}{\partial y} = -6y, \quad \frac{\partial z}{\partial y}\Big|_{(2,1)} = -6.
    \]
    \[
      \boxed{z = 4 + 12(x-2) - 6(y-1)}
    \]
  \end{solution}
  \vfill
  \vfill
  \vfill
  
\item Consider the function $f(x,y) = e^x \cos(y)$.
  \begin{enumerate}
  \item (6 points) Find the directional derivative of $f$ in
    the direction of the vector $\left< 3,-4 \right>$ at the point
    $(0,\frac{\pi}{2})$.
    \begin{solution}
      \[
        \nabla f = \vv{e^x\cos y}{-e^x \sin y}, \qquad
        \nabla f\big(0,\tfrac{\pi}{2}\big) = \vv{0}{-1}.
      \]
      $\left|\vv{3}{-4}\right| = 5$, so $\bu = \vv{\tfrac35}{-\tfrac45}$ and
      \[
        D_{\bu} f\big(0,\tfrac{\pi}{2}\big) = \vv{0}{-1}\cdot\vv{\tfrac35}{-\tfrac45}
        = 0 + \tfrac45 = \boxed{\frac45}
      \]
    \end{solution}
    \vfill
    \vfill
    
  \item (3 points) Find the direction of the maximum rate of change of $f$ at the
    point $(0,\frac{\pi}{2})$.
    \begin{solution}
      The maximum rate of change is in the direction of the gradient:
      \[
        \boxed{\nabla f\big(0,\tfrac{\pi}{2}\big) = \vv{0}{-1}}
      \]
    \end{solution}
    \vfill
    
  \item (3 points) Find the maximum rate of change of $f$ at the point
    $(0,\frac{\pi}{2})$.
    \begin{solution}
      $\left|\nabla f\big(0,\tfrac{\pi}{2}\big)\right| = \sqrt{0^2 + (-1)^2} = \boxed{1}$
    \end{solution}
    \vfill
    
  \end{enumerate}

  
  
\end{enumerate}

\end{document}
%%% Local Variables:
%%% mode: LaTeX
%%% TeX-master: t
%%% End:
